\documentclass[11pt]{article}
\usepackage[a4paper,margin=1in]{geometry}
\usepackage{amsmath,amssymb,amsthm,microtype}
\usepackage{cleveref}
\newtheorem{theorem}{Theorem}[section]
\newtheorem{corollary}[theorem]{Corollary}
\newtheorem{proposition}[theorem]{Proposition}
\theoremstyle{remark}
\newtheorem{remark}[theorem]{Remark}
\newcommand{\R}{\mathbb R}
\newcommand{\ka}{\kappa_{\mathrm{aff}}}
\newcommand{\TV}{\mathrm{TV}}
\title{An Oriented Paired Affine Gaussian Chain at Critical Clearance}
\author{Bin Wang}
\date{July 2026}
\begin{document}
\maketitle

\begin{abstract}
The critical endpoint row of RH-322 has a half-line Gaussian limit indexed by
the clearance phase.  We propagate that probability law through the two
signed affine tangent legs from the endpoint to the repelling boundary.  The
resulting joint law
\[
 J_d(v,u,w)=g_d(v)\phi(u+2u_cv)\phi(w+\lambda u)
\]
retains the source coordinate and therefore transfers the exact RH-322
$L^1$ error isometrically.  Its intermediate and output marginals are
explicit extended-skew-normal densities, not single Gaussians, and we give
their complete first and second moments.  At zero clearance, both nominal
orientations have positive leakage of order one.  This is an exact theorem
for the local affine probability chain.  It does not control the finite-noise
curvature and normalization remainder, combine the neighboring sibling or
parity layer, or prove a joint first-alias trace law.
\end{abstract}

\section{The signed tangent chain}

Let $u_c$ be the root in $(1,2)$ of
\[
 u^3-2u^2+2u-2=0,
 \qquad r=u_c-1,
 \qquad \lambda=2u_cr.
\]
The archived values are
\[
 u_c=1.543689012692076\ldots,
 \qquad \lambda=1.678573510428322\ldots.
\]
Set
\begin{equation}\label{eq:constants}
 \alpha=2u_c,
 \qquad \ka=\alpha\lambda,
 \qquad \beta=\sqrt{1+\lambda^2},
\end{equation}
so that
\[
 \alpha=3.087378025384153\ldots,
 \quad \ka=5.182390970088339\ldots,
 \quad \beta=1.953870269468181\ldots.
\]
The subscript in $\ka$ is essential: RH-14 already uses the unadorned
$\kappa$ for a different parity-profile constant.

The two local expansions extracted in RH-14 are
\begin{align}
 f(1-\sigma v)
  &=-r+\alpha\sigma v+O(\sigma^2v^2),
  \label{eq:endpoint-expansion}\\
 f(r+\sigma u)
  &=r-\lambda\sigma u-u_c\sigma^2u^2.
  \label{eq:repelling-expansion}
\end{align}
After folding the negative first image, the signed tangent coordinates obey
\begin{equation}\label{eq:oriented-chain}
 U=-\alpha V-Z_1,
 \qquad W=-\lambda U+Z_2,
\end{equation}
where $Z_1,Z_2$ are independent standard normals.  Hence
\begin{equation}\label{eq:collapsed-chain}
 W=\ka V+\lambda Z_1+Z_2
   =\ka V+\beta Z
\end{equation}
with $Z$ standard normal and independent of $V$.  The variable $W$ is
centered at the fixed point $r$.  RH-17 gives
$x_{k,1}-r=\ka\delta_k(1+o(1))$, so comparison with a packet centered at the
actual cycle point requires the translated coordinate $W-\ka d$.

Equations \eqref{eq:oriented-chain}--\eqref{eq:collapsed-chain} define the
limiting tangent model.  They are not asserted as exact finite-$\sigma$
physical dynamics; the omitted terms in
\eqref{eq:endpoint-expansion}--\eqref{eq:repelling-expansion}, fold changes,
and finite row normalizers are the RH-324 interface.

\section{Exact Markov lifting of the RH-322 row}

Write
\[
 \phi(t)=\frac{e^{-t^2/2}}{\sqrt{2\pi}},
 \qquad
 \Phi(t)=\int_{-\infty}^t\phi(s)\,ds,
 \qquad \overline\Phi=1-\Phi,
\]
and, for $d\ge0$,
\begin{equation}\label{eq:entrance-profile}
 g_d(v)=\frac{\phi(v-d)}{\Phi(d)}\mathbf1_{[0,\infty)}(v).
\end{equation}
For any probability density $p$ on $\R_+$, define its oriented affine lift
\begin{equation}\label{eq:affine-lift}
 (\mathcal A p)(v,u,w)
 =p(v)\phi(u+\alpha v)\phi(w+\lambda u).
\end{equation}

\begin{theorem}[Joint affine transfer isometry]\label{thm:isometry}
For any probability densities $p,q$ on $\R_+$,
\begin{equation}\label{eq:isometry}
 \|\mathcal A p-\mathcal A q\|_{L^1(\R_+\times\R^2)}
 =\|p-q\|_{L^1(\R_+)}.
\end{equation}
Every marginalization of the lifted laws is an $L^1$ contraction.
\end{theorem}

\begin{proof}
The two conditional Gaussian densities in \eqref{eq:affine-lift} are
nonnegative and each integrates to one.  Tonelli's theorem gives
\[
 \int_{\R_+\times\R^2}
 |p(v)-q(v)|\phi(u+\alpha v)\phi(w+\lambda u)\,dv\,du\,dw
 =\int_{\R_+}|p-q|.
\]
Integrating out coordinates before taking the absolute value gives the
marginal contraction statement.
\end{proof}

To retain the exact RH-322 physical seed, put $L=\sigma^{-1}$ and
\begin{equation}\label{eq:finite-seed}
 \widetilde h_{\sigma,a}(v)=
 \frac{\phi(v-a)+\phi(2L-v-a)}
 {\Phi(a)-\overline\Phi(2L-a)}\mathbf1_{[0,L]}(v).
\end{equation}
Define
\[
 J_{\sigma,a}=\mathcal A\widetilde h_{\sigma,a},
 \qquad J_d=\mathcal A g_d.
\]
The first object is an exact physical endpoint row followed by the exact
\emph{affine} lift; it is not the actual finite-noise two-step row.

\begin{corollary}[Finite-seed joint error]\label{cor:finite-seed}
For $\sigma>0$, $0\le a\le\sigma^{-1}$, and $d\ge0$,
\begin{align}
 \|J_{\sigma,a}-J_d\|_1
 &=\|\widetilde h_{\sigma,a}-g_d\|_1,
 \label{eq:joint-base-equality}\\
 \|J_{\sigma,a}-J_d\|_1
 &\le |a-d|+
 \frac{2\overline\Phi(\sigma^{-1}-a)}{\Phi(a)}.
 \label{eq:joint-bound}
\end{align}
In particular,
\begin{equation}\label{eq:same-phase}
 \|J_{\sigma,a}-J_a\|_1
 =\frac{2\overline\Phi(\sigma^{-1}-a)}{\Phi(a)}.
\end{equation}
The same right side bounds every marginal distance.  Under the convention
$d_{\TV}=\tfrac12\|\cdot\|_1$, all statements divide by two.
\end{corollary}

\begin{proof}
The identity is \cref{thm:isometry}.  RH-322 proves the exact same-phase
tail in \eqref{eq:same-phase} and
$\|g_a-g_d\|_1\le|a-d|$; the triangle inequality gives
\eqref{eq:joint-bound}.
\end{proof}

The equality in \eqref{eq:joint-base-equality} requires retaining $V$.
After $V$ is marginalized, only contraction is available; no marginal
isometry is claimed.

\section{Explicit marginals and a non-Gaussian output}

Put
\begin{equation}\label{eq:scales}
 s_1=\sqrt{1+\alpha^2},
 \qquad s_2=\sqrt{\ka^2+\beta^2}.
\end{equation}

\begin{proposition}[Selected-Gaussian marginals]\label{prop:marginals}
The intermediate and output densities of $J_d$ are
\begin{align}
 p_d(u)
 &=\frac1{s_1\Phi(d)}
 \phi\!\left(\frac{u+\alpha d}{s_1}\right)
 \Phi\!\left(\frac{d-\alpha u}{s_1}\right),
 \label{eq:intermediate-density}\\
 q_d(w)
 &=\frac1{s_2\Phi(d)}
 \phi\!\left(\frac{w-\ka d}{s_2}\right)
 \Phi\!\left(
 \frac{\ka w+\beta^2d}{\beta s_2}
 \right).
 \label{eq:output-density}
\end{align}
The joint $(U,W)$ density is $p_d(u)\phi(w+\lambda u)$.
\end{proposition}

\begin{proof}
Let $V_0\sim N(d,1)$ before conditioning on $V_0\ge0$.  The unconditioned
$U_0=-\alpha V_0-Z_1$ has law $N(-\alpha d,s_1^2)$, while
\[
 V_0\mid U_0=u
 \sim N\!\left(\frac{d-\alpha u}{s_1^2},\frac1{s_1^2}\right).
\]
Multiplying the unconditioned $U_0$ density by
$\Pr(V_0\ge0\mid U_0=u)/\Phi(d)$ proves
\eqref{eq:intermediate-density}.  Likewise
$W_0=\ka V_0+\beta Z$ has law $N(\ka d,s_2^2)$ and
\[
 V_0\mid W_0=w
 \sim N\!\left(
 \frac{\ka w+\beta^2d}{s_2^2},
 \frac{\beta^2}{s_2^2}
 \right),
\]
which gives \eqref{eq:output-density}.  The joint formula follows directly
from $W\mid U=u\sim N(-\lambda u,1)$.
\end{proof}

The output is not a single Gaussian.  Indeed, let
\[
 h_d(w)=\frac1{s_2}\phi\!\left(\frac{w-\ka d}{s_2}\right).
\]
Then \eqref{eq:output-density} gives the exact ratio
\begin{equation}\label{eq:tail-ratio}
 \frac{q_d(w)}{h_d(w)}
 =\frac1{\Phi(d)}
 \Phi\!\left(\frac{\ka w+\beta^2d}{\beta s_2}\right)
 \longrightarrow\frac1{\Phi(d)}>1
 \quad(w\to+\infty)
\end{equation}
for every finite $d\ge0$.  A ratio of two Gaussian densities can have a
finite nonzero right-tail limit only when their means and variances agree;
then the ratio is identically one.  Thus \eqref{eq:tail-ratio} proves the
scoped negative result that $q_d$ is not any ordinary Gaussian.

\section{Moments and surviving orientations}

Let
\begin{equation}\label{eq:mills}
 r_d=\frac{\phi(d)}{\Phi(d)},
 \qquad m_d=d+r_d,
 \qquad \nu_d=1-dr_d-r_d^2.
\end{equation}

\begin{proposition}[Mean and covariance]\label{prop:moments}
The mean vector is
\begin{equation}\label{eq:mean-vector}
 \mathbb E(V,U,W)=m_d(1,-\alpha,\ka).
\end{equation}
The nonredundant covariance entries are
\begin{align}
 \operatorname{Var}V&=\nu_d,&
 \operatorname{Var}U&=\alpha^2\nu_d+1,&
 \operatorname{Var}W&=\ka^2\nu_d+\beta^2,
 \label{eq:variances}\\
 \operatorname{Cov}(V,U)&=-\alpha\nu_d,&
 \operatorname{Cov}(V,W)&=\ka\nu_d,&
 \operatorname{Cov}(U,W)&=-\lambda(\alpha^2\nu_d+1).
 \label{eq:covariances}
\end{align}
Relative to the deterministic centers $(d,-\alpha d,\ka d)$, the
conditioning bias is
\begin{equation}\label{eq:bias}
 r_d(1,-\alpha,\ka).
\end{equation}
\end{proposition}

\begin{proof}
The RH-322 entrance law has mean $m_d$ and variance $\nu_d$.  Apply the
independent-noise identities \eqref{eq:oriented-chain} and
\eqref{eq:collapsed-chain}.  In particular,
$\operatorname{Cov}(U,W)=-\lambda\operatorname{Var}(U)$; the first-leg noise
prevents replacing this entry by a deterministic-product formula.
\end{proof}

At $d=0$, planar Gaussian wedge probabilities give
\begin{align}
 \Pr(U>0)
 &=\frac1\pi\arctan\frac1\alpha
 =0.0997061646699731\ldots,
 \label{eq:u-leak}\\
 \Pr(W<0)
 &=\frac1\pi\arctan\frac\beta\ka
 =0.1147638712758368\ldots.
 \label{eq:w-leak}
\end{align}
These are unconditional marginal probabilities.  They show that the
nominally negative $U$ and positive $W$ orientations both have surviving
opposite-sign mass in the affine limit.  They are not parity weights and do
not determine a conditional branch-return probability.

\section{Protocol and route firewall}

The reproduction code uses only the standard library.  Composite Simpson
quadrature checks the two closed-form marginals against their defining
convolutions, their normalization and moments, the exact finite-seed tail,
and the two wedge probabilities.  No parameter fitting or finite-order
spectral inference enters the argument.

This probability calculation must also be separated from RH-18.  That paper
propagates peak-normalized Gaussian observables by a backward affine operator
and obtains an amplitude coefficient from peak normalization.  Here every
object is a forward probability density of mass one; no RH-18 packet
coefficient is a probability-loss factor.  RH-19 further proves that the
neighboring critical sibling has order-one natural $L^2$ mass.  The sibling
is not contained in $J_d$ and cannot be dismissed as a Gaussian tail.

\begin{remark}[Claim boundary]
No actual finite-noise two-leg curvature or normalization remainder is
proved.  No parity cancellation, sibling/shell matching, moving-order
Duhamel composition, branch-isolated return, full-trace replacement, or
joint error $o(kR^{-2k})$ with $R=1.4$ is obtained.  Gates A--E remain
false/open.  The paper constructs no Hilbert--Polya operator, identifies no
Riemann zero, proves no von Mangoldt trace formula or zeta-divisor equality,
and does not imply RH.
\end{remark}

The next interface is RH-324: turn one tangent leg into a quantitative
finite-noise physical-kernel approximation by controlling curvature, fold
changes, state-boundary tails, and exact row normalization on one weighted
domain.  Moving-order composition begins only after that local remainder is
available.

\end{document}
